\documentclass[9pt,a4paper,twocolumn]{ltjsarticle}
\usepackage[top=20mm,bottom=22mm,left=17mm,right=17mm,columnsep=8mm]{geometry}
\usepackage{amsmath}
\usepackage{booktabs}
\usepackage[hidelinks]{hyperref}
\newcommand{\NSamples}{442}
\newcommand{\NFeatures}{10}
\newcommand{\NTrain}{353}
\newcommand{\NTest}{89}
\newcommand{\Seed}{0}
\newcommand{\TargetMean}{152.1}
\newcommand{\TargetStd}{77.0}
\newcommand{\NPolyFeatures}{65}
\newcommand{\TrainMse}{2256.1}
\newcommand{\TrainRmse}{47.5}
\newcommand{\TrainMae}{38.0}
\newcommand{\TrainRsq}{0.632}
\newcommand{\TestMse}{3637.1}
\newcommand{\TestRmse}{60.3}
\newcommand{\TestMae}{46.4}
\newcommand{\TestRsq}{0.291}
\newcommand{\CpuModel}{AMD EPYC 9V45 96-Core Processor}
\newcommand{\CpuCores}{4}
\newcommand{\MemoryGib}{15.6}
\newcommand{\PythonVersion}{3.13.3}
\newcommand{\NumpyVersion}{2.3.4}
\newcommand{\SklearnVersion}{1.7.2}

% 各文の ID。claims.json の sentences の ID と同じ
% 各文の ID。この文の最初の claims の行（本リポジトリの claims.json、commit 93295a4）へリンクする
\newcommand{\sid}[1]{\textsuperscript{\textsf{\scriptsize[\href{https://github.com/auto-res2/diabetes-polynomial-regression-4/blob/93295a4/formal/claims.json\#L\csname line#1\endcsname}{#1}]}}}
\expandafter\def\csname lineR01\endcsname{86}
\expandafter\def\csname lineR02\endcsname{373}
\expandafter\def\csname lineR03\endcsname{458}
\expandafter\def\csname lineR04\endcsname{546}
\expandafter\def\csname lineR05\endcsname{557}
\expandafter\def\csname lineR06\endcsname{583}
\expandafter\def\csname lineR07\endcsname{652}
% 論文の途中からの抜粋: 節と頁の番号を途中から始める
\setcounter{section}{3}
\setcounter{page}{3}
\setcounter{equation}{2}
\setcounter{table}{1}
\setcounter{footnote}{1}

% 文献 [1] の claims.json の、引用した claims の行へのリンク
\newcommand{\clm}[2]{\href{https://github.com/auto-res2/diabetes-linear-regression-6/blob/bc1e74f/formal/claims.json\#L#2}{#1}}
\begin{document}

\twocolumn[{
\begin{center}
{\Large\bfseries 糖尿病データセットに対する多項式回帰の予測性能}\\[2mm]
{\normalsize auto-res2}\\[1mm]
{\small （「4 実験設定」および「5 結果」の抜粋）}
\end{center}
\vspace{3mm}
}]

\section{実験設定}

\sid{R01}データセット，前処理，訓練データとテストデータへの分割は，文献~\cite{report1} と同じとした~\cite[\clm{R01}{32}, \clm{R04}{323}, \clm{R06}{453}]{report1}．\sid{R02}予測性能は，訓練データとテストデータのそれぞれについて，MSE，RMSE，MAE，$R^2$ で評価した~\cite[\clm{R11}{709}, \clm{R12}{742}, \clm{R13}{762}, \clm{R14}{782}, \clm{R15}{897}]{report1}．\sid{R03}予測器には，説明変数を 2 次の多項式（交互作用項を含む計 \NPolyFeatures{} 個の変数）に展開し，訓練データに対して最小二乗法で当てはめたモデルを用いた．\sid{R04}正則化は行っていない．

\section{結果}

\sid{R05}テストデータに対する MSE は \TestMse{} であった．\sid{R06}これは文献~\cite{report1} の線形回帰のテストデータに対する MSE 3424.3~\cite[\clm{R20}{1049}]{report1} より小さく，多項式フィッティングは線形回帰より予測精度が高い．

\begin{table}[t]
\centering
\caption{\sid{R07}予測器ごとの評価指標．線形回帰の値は文献~\cite{report1} の表 2~\cite[\clm{R26}{1620}]{report1}．}
\label{tab:metrics}
\begin{tabular}{llrrrr}
\toprule
予測器 & データ & MSE & RMSE & MAE & $R^2$ \\
\midrule
線形回帰~\cite{report1} & 訓練 & 2734.8 & 52.3 & 42.6 & 0.554 \\
 & テスト & 3424.3 & 58.5 & 46.2 & 0.332 \\
2 次の多項式 & 訓練 & \TrainMse & \TrainRmse & \TrainMae & \TrainRsq \\
 & テスト & \TestMse & \TestRmse & \TestMae & \TestRsq \\
\bottomrule
\end{tabular}
\end{table}

\begin{thebibliography}{9}
\small
\bibitem{report1} auto-res2: 糖尿病データセットに対する線形回帰モデルの予測性能, \url{https://github.com/auto-res2/diabetes-linear-regression-6}, commit bc1e74f.
\end{thebibliography}

\end{document}
